% =========================================================================
% REFLECTION ON LEADERSHIP & PERSONAL GROWTH
% =========================================================================

\chapter{Reflection on Leadership \& Personal Growth}
\label{app:reflection}

\section{Leadership Journey}
\label{sec:journey}

When I began this thesis, I thought leadership meant leading others. Now that I have gotten to the end of it, I understood that it also means leading yourself through the parts of an innovative project that probably no one else has ever done before and can do for you. My experience over these months was a mix of the two, and I have come to value both.

At the beginning of this thesis, I went through several iterations of the research question, partly because I was still figuring out what I truly wanted to do. I knew I wanted to create something innovative, meaningful, and capable of having a real impact, while bringing together the skills I had developed in recent years in a creative and unconventional way.

This reflects how I approach things: I love learning, exploring, connecting ideas across fields, and creating something both rigorous and meaningful. The thesis gave me the opportunity to do this, but also challenged me to work through uncertainty and complexity.

I knew the topic could easily become too broad or partially drift away, so narrowing it down was a major challenge. I had to constantly separate what was simply interesting from what was truly relevant to the thesis. I believe I managed to find that balance while keeping the work focused and coherent.

Building the analysis, designing the visualisations, running the study, coding the responses, and writing the chapters all required me to make decisions and take responsibility for them. I learned that leadership is often quiet: accepting results that do not support what I expected, and not making the data say more than it can.

More than once, I had to accept that the findings were different from what I had hoped for and explain why they were still meaningful. Being honest about this was one of the most important leadership lessons I learned.

I also learned to lead a project through change. The study I completed is not the one I originally imagined. Some ideas had to be reduced, methods reconsidered, and ambitions left for future work. I have to admit that this, at first, felt a bit like failure. In the end, I understood it as good judgement: knowing what to let go of so that what remains can be done well and I am glad I made it this far.


\section{Inspirational Leadership}
\label{sec:inspiration}

The people who inspire me are rarely the ones who do a single thing perfectly. They are the eclectic ones: the systemic, transversal, curious minds who hold real expertise but refuse to be confined by it, who move across fields and let their many passions feed one another. I am drawn to people who never stop learning and, in doing so, make others want to learn too.

Federico Faggin is, for me, the clearest example. He designed the first microprocessor, the chip at the heart of nearly every device we now use, and then, in what he calls his later life, turned that same rigorous mind toward the study of consciousness, seeking to bring physics and the inner life of experience into one picture. What moves me is not any single claim of his, but the shape of his path: a person willing to be a beginner again, to carry the discipline of one domain into a completely different one, and to ask the largest questions without abandoning precision. He gave me permission, in a sense, to approach a problem differently, to trust that a technical tool and a human question belong in the same sentence.

I feel the same pull toward others who cross boundaries. Carlo Rovelli, who writes about physics as if it were poetry and never loses the science in the beauty. Michela Murgia, who used voice and presence to widen how people see. Jane Goodall, who changed a field by paying a new kind of attention, patiently, for years, and then spent her life communicating what she had seen. I am aware these can sound like different domains entirely, but I can guarantee that the underlying quality is the same: a refusal to separate deep understanding from the will to share it.

That quality is, quietly, part of what this thesis is about. The idea of \emph{enhanced cognition}, using technology not to replace human understanding but to expand it, is present throughout the whole project. I like to think of visualisation as a way of extending what we are able to perceive, allowing us to see and understand parts of the world that would otherwise remain abstract or invisible. For me, this thesis is a small technical example of a much bigger belief: that when technology is connected with human understanding, it can help us see the world differently, feel more connected to it, and perhaps also become more aware of our responsibility towards it.

\section{Communicating Impactfully}

\label{sec:impact}

If there is one idea that connects my values with my work, it is the belief that understanding is only complete when it can be shared. In this sense, the whole thesis is also about communication: taking a system that is almost impossible to see directly and finding a way to make it understandable, without simplifying it to the point of losing what makes it real.

Writing this thesis taught me a lot about communicating with impact. I learned that clarity does not mean losing rigour. On the contrary, a precise idea should be expressed in clear and human language. I also learned that the reader's attention cannot be taken for granted, and that communicating well means taking responsibility for how a message is received. I had to explain things such as a machine-learning pipeline, a statistical null result, or a participant's misconception in ways that could also be understood by someone outside the field.

The user study made this lesson even clearer. Seeing non-specialists describe what they understood from the visualisations, sometimes correctly and sometimes with great confidence but in the wrong way, reminded me that communication is not only about what we intend to say. It is about what the other person actually understands. This made me a more careful and honest communicator, and showed me that impact is not measured by how much we say, but by what reaches the other person.

\section{Alignment with Career Vectors}

\label{sec:alignment}

Looking forward, this work fits closely with the direction I want to take. I have always been interested in the space between science, technology, creativity, and human understanding, and I can imagine continuing towards physics or another interdisciplinary field. This thesis was, in many ways, a first experience of that direction. It allowed me to combine analytical work, design, empirical research, and communication within one project, and confirmed that this is the kind of work that motivates me most.

I do not see myself working only within one narrow field, but I also do not want to be a generalist without depth. I want to build strong expertise while remaining curious and able to move across disciplines. I can see myself working at the intersection of technology and creativity, in an innovative environment where I can use different skills to turn complex ideas into something meaningful and useful. At the same time, I want to continue studying and learning, especially in science and technology, to build a broader and deeper understanding of the world rather than stopping at one field.

This direction is closely connected to the mission that has guided me throughout this work: \emph{Wisdom Unlocks Awareness, and Awareness Unlocks Agency}. For me, knowledge becomes meaningful when it creates awareness, and awareness becomes meaningful when it enables people to understand, decide, and act. This thesis has been one small attempt to put that idea into practice: using technology and creativity to make something difficult to see easier to understand.

If this thesis has taught me something beyond its results, it is that I can take a complex and uncertain idea, develop it through many different stages, and bring it to a finished whole while staying honest about its limits. That is probably what I will carry with me most. Not one single result, but the confidence that I can keep exploring, keep learning, and connect different forms of knowledge in order to create something that matters.